\documentclass[11pt,oneside]{article}
\usepackage{subfigure}
\usepackage{amsmath}
\usepackage{amssymb}
\usepackage{amsbsy}
\usepackage{mathpazo}
\usepackage{float}
\usepackage{braket}
\setlength {\marginparwidth }{3cm}
\usepackage{todonotes}
\usepackage[spanish]{babel}
\usepackage{graphicx}
\usepackage{listings}
\usepackage{color}
\usepackage{booktabs}
\usepackage{multirow}
\usepackage{ragged2e}
\usepackage{amsmath,blkarray,booktabs,bigstrut}
\usepackage[a4paper, left=2.5cm, right=2.5cm, top=2.5cm, bottom=2.5cm]{geometry}
\usepackage{hyperref}
\usepackage{colortbl}
\usepackage{xcolor}
\usepackage{enumitem}
\usepackage{tcolorbox}
\usepackage{todonotes}

\hypersetup{
    colorlinks=true,
    linkcolor=blue, % Color de los enlaces internos
    filecolor=magenta, % Color de los enlaces a archivos
    urlcolor=blue % Color de los enlaces a URLs
}

\floatstyle{ruled}
\restylefloat{table}

\usepackage{listings}
\usepackage{color}
\lstloadlanguages{csh}

\definecolor{red}{rgb}{0.6,0,0} 
\definecolor{blue}{rgb}{0,0,0.6}
\definecolor{green}{rgb}{0,0.8,0}
\definecolor{cyan}{rgb}{0.0,0.6,0.6}
\definecolor{white}{rgb}{1,1,1}

% Configuración global para todos los recuadros
\tcbset{
    colback=blue!3!white,     % Fondo
    colframe=blue!50!white,   % Borde 
    boxrule=0.5mm,            % Grosor del borde (opcionalmente más delgado)
    title=Información         % Título predeterminado
}

\lstset{
language=csh,
captionpos=b,      % Posición del caption (b: abajo, t: arriba)
basicstyle=\footnotesize\ttfamily,
% numbers=left,
% numberstyle=\tiny,
% numbersep=5pt,
tabsize=2,
extendedchars=true,
% breaklines=true,
% frame=b,
stringstyle=\color{red}\ttfamily,
showspaces=false,
showtabs=false,
% xleftmargin=17pt,
framexleftmargin=17pt,
% framexrightmargin=5pt,
framexbottommargin=4pt,
commentstyle=\color{green},
morecomment=[l]{//}, %use comment-line-style!
morecomment=[s]{/*}{*/}, %for multiline comments
showstringspaces=false,
morekeywords={ abstract, event, new, struct,
as, explicit, null, switch,
base, extern, object, this,
bool, false, operator, throw,
break, finally, out, true,
byte, fixed, override, try,
case, float, params, typeof,
catch, for, private, uint,
char, foreach, protected, ulong,
checked, goto, public, unchecked,
class, if, readonly, unsafe,
const, implicit, ref, ushort,
continue, in, return, using,
decimal, int, sbyte, virtual,
default, interface, sealed, volatile,
delegate, internal, short, void,
do, is, sizeof, while,
double, lock, stackalloc,
else, long, static, get,
enum, namespace, string},
keywordstyle=\color{blue},
identifierstyle=\color{black},
backgroundcolor=\color{white},
inputencoding=utf8,
    literate=%
    {á}{{\'a}}1
    {é}{{\'e}}1
    {í}{{\'i}}1
    {ó}{{\'o}}1
    {ú}{{\'u}}1
    {Á}{{\'A}}1
    {É}{{\'E}}1
    {Í}{{\'I}}1
    {Ó}{{\'O}}1
    {Ú}{{\'U}}1
    {ñ}{{\~n}}1
    {Ñ}{{\~N}}1
    {¡}{{\textexclamdown}}1
    {¿}{{\textquestiondown}}1
}

% Definir tipos personalizados en cian
\lstset{
    classoffset=1,
    morekeywords={
    Program, Console, List, StringBuilder, Array,
    TypeOfTriangle, DayOfWeek,
    NotImplementedException, DivideByZeroException, ArgumentException, 
    ChessPiece,
    BigInt, Fraction, Poly, Set, Book, Genre},
    keywordstyle=\color{cyan},
    classoffset=0,
}

\renewcommand{\lstlistingname}{Figura} % Cambia el prefijo de lstlisting a "Figura"

% Define la variable para controlar la visibilidad de las respuestas
\newif\ifshowanswers

% Define la variable que representa el curso
\def\academicyear{2024 - 2025}

% \showanswerstrue  % Comenta esta línea para ocultar las respuestas

\begin{document}
    \begin{center}
    \begin{large}
    Cp 9 - Tableros\\
    Curso \academicyear\\
    \end{large}
    \begin{figure}[h]
    	\centering
    	\includegraphics[width=0.5\linewidth]{cp7/matrix.png}
    \end{figure}
\end{center}

\section{4 en línea}
\begin{enumerate}[label=\alph*)]
    \item Dada una matriz de elementos booleanos, implemente un método que determine si hay 4 valores idénticos en línea. Estos valores pueden estar alineados en dirección horizontal, vertical o diagonal.
    \item * Lo mismo que el inciso anterior pero en lugar de devolver si existe o no, devolver cuántas líneas de 4 hay.
\end{enumerate}


\textbf{Ejemplo:}

Supongamos la siguiente matriz donde los valores \texttt{true} están marcados con una ''T'' y los valores \texttt{false} con una ''F'':

\[
\begin{array}{cccccc}
F & F & T & F & T & F \\
\textcolor{green}{T} & \textcolor{green}{T}& \textcolor{green}{T} & \textcolor{green}{T} & F & \textcolor{red}{F} \\
F & \textcolor{green}{T} & F & T & \textcolor{red}{F} & T \\
T & F & \textcolor{green}{T} & \textcolor{red}{F} & T & F \\
T & T & \textcolor{red}{F} & \textcolor{green}{T} & F & F \\
\end{array}
\]

En este ejemplo, existen tres posibles líneas de 4 valores iguales consecutivos.


\section{Tic-Tac-Toe}
Implementa un juego sencillo de Cerito-Cruz para dos jugadores humanos. Ten en cuenta los siguientes requisitos:
\begin{itemize}
    \item Usa un array bidimensional de 3x3 para representar el tablero.
    \item Los jugadores deben turnarse para colocar sus símbolos (\texttt{X} o \texttt{O}). Comprueba que los movimientos sean válidos.
    \item En cada turno verifica si algún jugador ha ganado o si hay un empate.
\end{itemize}

\section{Sopa de letras}
Implemente el método \texttt{bool EstaEnLaSopa(char[,] sopa, string palabra)} que determine si el string \texttt{palabra} se encuentra en \texttt{sopa}. 
\begin{itemize}
    \item Ejemplo: 
    \begin{lstlisting}
    string palabra = "pelo"; 
    char[,] sopa = { 
        { 'a', 'a', 'f', 't', 'j', 'q', 'w', 'e', 'r', 'o', 'p' }, 
        { 'g', 'j', 'p', 'b', 'j', 'e', 'r', 'o', 'a', 's', 'k' }, 
        { 'l', 'x', 'c', 'e', 't', 'y', 'e', 'r', 'a', 'o', 'n' }, 
        { 'b', 'g', 'j', 'f', 'l', 'd', 'e', 'r', 's', 't', 'o' }, 
        { 'q', 'u', 'e', 'r', 't', 'o', 'g', 's', 'e', 'm', 't' } 
    }; 
    EstaEnLaSopa(sopa, palabra); // true
    \end{lstlisting}
    \textbf{Nota:} Solo palabras que se lean de izquierda a derecha, de arriba hacia abajo y la diagonalcorrespondiente.
\end{itemize}


\section{Tablero alcanzable por caballos}
Dado un tablero de ajedrez donde se ubican caballos, determine si el tablero es completamente alcanzable (todas sus casillas son alcanzables). Una casilla es alcanzable si algún caballo puede llegar a ella en 0 o 1 movimientos. La representación del tablero será un array bidimensional de bool donde solo habrá true en las casillas donde haya un caballo. 
\begin{itemize}
    \item En este caso se devuelve false ya que solo se alcanzan las casillas \( (0,3), (1,1), (2,3),(3,0), (3,2) \).
    \begin{center}
    \begin{tikzpicture}[scale=0.8]
        % Dibuja el tablero
        \foreach \x in {0,1,2,3} {
            \foreach \y in {0,1,2,3} {
                \fill[white] (\x,\y) rectangle (\x+1,\y+1);
                \draw (\x,\y) rectangle (\x+1,\y+1);
            }
        }
    
        % Marcar los caballos con C_i
        \node at (1.5, 1.5) {$C_1$}; % (1,1)
    
        % Marcar las casillas alcanzables con X_i
        \node at (0.5, 3.5) {$X_1$}; % (0,3)
        \node at (1.5, 1.5) {$C_1$}; % (1,1)
        \node at (2.5, 3.5) {$X_1$}; % (2,3)
        \node at (3.5, 0.5) {$X_1$}; % (3,0)
        \node at (3.5, 2.5) {$X_1$}; % (3,2)
    
        % Etiquetas
        \foreach \x in {0,...,3} {
            \foreach \y in {0,...,3} {
                % Etiquetas de fila y columna
                \ifnum\x=0
                    \node at (-0.5,\y+0.5) {\textbf{\small \y}};
                \fi
                \ifnum\y=0
                    \node at (\x+0.5,-0.5) {\textbf{\small \x}};
                \fi
            }
        }
    \end{tikzpicture}
    \end{center}
    \item En otro caso se devuelve true pues todas las casillas son alcanzables (algunas por más de un caballo).
    \begin{center}
    \begin{tikzpicture}[scale=0.8]
        % Dibuja el tablero
        \foreach \x in {0,1,2,3} {
            \foreach \y in {0,1,2,3} {
                \fill[white] (\x,\y) rectangle (\x+1,\y+1);
                \draw (\x,\y) rectangle (\x+1,\y+1);
            }
        }
        
        % Caballos en posiciones C_i
        \node at (1.5, 1.5) {$C_1$}; % (1,1)
        \node at (2.5, 1.5) {$C_2$}; % (2,1)
        \node at (1.5, 2.5) {$C_3$}; % (1,2)
        \node at (2.5, 2.5) {$C_4$}; % (2,2)
    
        % Marcar las casillas alcanzables con X_i
        \node at (0.5, 0.5) {$X_2$}; % (0,0)
        \node at (0.5, 1.5) {$X_4$}; % (0,1)
        \node at (0.5, 2.5) {$X_2$}; % (0,2)
        \node at (0.5, 3.5) {$X_1$}; % (0,3)
        \node at (1.5, 0.5) {$X_4$}; % (1,0)
        \node at (1.5, 3.5) {$X_2$}; % (1,3)
        \node at (2.5, 0.5) {$X_3$}; % (2,0)
        \node at (2.5, 3.5) {$X_1$}; % (2,3)
        \node at (3.5, 0.5) {$X_1$}; % (3,0)
        \node at (3.5, 1.5) {$X_3$}; % (3,1)
        \node at (3.5, 2.5) {$X_1$}; % (3,2)
        \node at (3.5, 3.5) {$X_2$}; % (3,3)
    
        % Etiquetas
        \foreach \x in {0,...,3} {
            \foreach \y in {0,...,3} {
                % Etiquetas de fila y columna
                \ifnum\x=0
                    \node at (-0.5,\y+0.5) {\textbf{\small \y}};
                \fi
                \ifnum\y=0
                    \node at (\x+0.5,-0.5) {\textbf{\small \x}};
                \fi
            }
        }
        
    \end{tikzpicture}
    \end{center}
\end{itemize}

\section{Dama amenazada}
Dado un tablero de ajedrez donde se ubican solo damas (reinas), determine si existe alguna que amenace a otra. La representación del tablero será un array bidimensional de bool donde solo habrá true en las casillas donde haya una reina.


\section{Validación de un Sudoku}
En este ejercicio, implementarás un programa que valide si un tablero de Sudoku es válido. El programa recibirá como entrada un tablero representado como una matriz de \(9 \times 9\), un tablero de Sudoku debe cumplir con las siguientes reglas:
\begin{itemize}
    \item Cada fila debe contener los números del \(1\) al \(9\) sin repetir.
    \item Cada columna debe contener los números del \(1\) al \(9\) sin repetir.
    \item Cada subcuadrícula de \(3 \times 3\) debe contener los números del \(1\) al \(9\) sin repetir.
\end{itemize}

Algunas celdas estarán vacías, representadas por un \(0\).

\subsection*{Ejemplo de entrada}
\[
\text{sudoku} = 
\begin{bmatrix}
5 & 3 & 0 & 0 & 7 & 0 & 0 & 0 & 0 \\
6 & 0 & 0 & 1 & 9 & 5 & 0 & 0 & 0 \\
0 & 9 & 8 & 0 & 0 & 0 & 0 & 6 & 0 \\
8 & 0 & 0 & 0 & 6 & 0 & 0 & 0 & 3 \\
4 & 0 & 0 & 8 & 0 & 3 & 0 & 0 & 1 \\
7 & 0 & 0 & 0 & 2 & 0 & 0 & 0 & 6 \\
0 & 6 & 0 & 0 & 0 & 0 & 2 & 8 & 0 \\
0 & 0 & 0 & 4 & 1 & 9 & 0 & 0 & 5 \\
0 & 0 & 0 & 0 & 8 & 0 & 0 & 7 & 9 \\
\end{bmatrix}
\]
\subsection*{Salida esperada}

\begin{itemize}
    \item Para el tablero de entrada anterior, el programa debe devolver:
    \[
    \text{El tablero de Sudoku es válido: Sí.}
    \]
    \item Si se modifica la celda \((1,2)\) del tablero para que contenga un \(5\), el programa debe devolver:
    \[
    \text{El tablero de Sudoku es válido: No.}
    \]
\end{itemize}

\section{Validación de movimientos en Ajedrez}
Tu misión es implementar un programa que determine si un movimiento de una pieza de ajedrez es válido. Usa el siguiente \textcolor{blue}{enum} para representar las piezas:

\begin{lstlisting}
public enum ChessPiece
{
    None,           // No piece (empty square)
    WhitePawn,
    WhiteKnight,
    WhiteBishop,
    WhiteRook,
    WhiteQueen,
    WhiteKing,
    BlackPawn,
    BlackKnight,
    BlackBishop,
    BlackRook,
    BlackQueen,
    BlackKing
}
\end{lstlisting}

El tablero se representará como una matriz de \(8 \times 8\) de tipo \texttt{ChessPiece}, donde cada celda indica la pieza presente o está vacía (\texttt{None}).

Tu programa debe recibir como entrada un tablero, junto con las coordenadas iniciales y finales del movimiento, representadas como \((x_1, y_1)\) y \((x_2, y_2)\); a partir de esta información, deberá determinar si el movimiento indicado es válido, considerando las reglas correspondientes a la pieza ubicada en la posición inicial.

Las reglas de movimiento para cada tipo de pieza son las siguientes:
\begin{itemize}
    \item \textbf{Peones} (Pawn): Avanzan (o retroceden, si son negros) una casilla hacia adelante, excepto en su primer movimiento, donde pueden avanzar dos casillas. Capturan piezas únicamente en diagonal y no pueden avanzar hacia casillas ocupadas.
    \item \textbf{Caballos} (Knight): Se mueven en forma de "L", moviéndose dos casillas en una dirección y una en perpendicular. Pueden saltar sobre otras piezas, pero solo ocupan la casilla de destino si está vacía o contiene una pieza del oponente.
    \item \textbf{Alfiles} (Bishop): Se mueven únicamente en diagonal, cualquier cantidad de casillas. No pueden atravesar ni ocupar casillas ocupadas, salvo para capturar una pieza del oponente en la casilla de destino.
    \item \textbf{Torres} (Rook): Se mueven en línea recta, ya sea horizontal o verticalmente, cualquier número de casillas. No pueden atravesar ni ocupar casillas ocupadas, salvo para capturar una pieza del oponente en la casilla de destino.
    \item \textbf{Reinas} (Queen): Combinan los movimientos de las torres y los alfiles, moviéndose tanto en línea recta como en diagonal.
    \item \textbf{Reyes} (King): Se mueven una casilla en cualquier dirección: horizontal, vertical o diagonal. No pueden moverse a una casilla ocupada salvo para capturar una pieza del oponente en la casilla de destino.
\end{itemize}

\textbf{Nota}: El ejercicio no requiere implementar lógica adicional como detectar jaque, jaque mate, enroque o promoción de peones. Solo se debe comprobar si el movimiento en cuestión cumple con las reglas básicas de la pieza.

\section{Jaque al rey}
Implemente un método que reciba un tablero de ajedrez, y determine si el rey negro está amenazado por alguna pieza blanca. 

El método deberá tener la siguiente signatura:

\begin{lstlisting}
    public enum ChessPiece
    {
        None,           // No piece (empty square)
        WhitePawn,
        WhiteKnight,
        WhiteBishop,
        WhiteRook,
        WhiteQueen,
        WhiteKing,
        BlackPawn,
        BlackKnight,
        BlackBishop,
        BlackRook,
        BlackQueen,
        BlackKing
    }
    
    public static bool IsBlackKingInCheck(ChessPiece[,] board)
    {
        // ...
    }
\end{lstlisting}


\section{El Mago y el Bosque Encantado}
En un bosque encantado, un mago ha decidido medir el poder de su hechizo de teletransportación. El bosque está representado como una cuadrícula de \(N \times M\), donde cada celda puede estar despejada o bloqueada por un obstáculo mágico. El mago comienza en una posición definida por el usuario, y desea calcular cuán lejos puede llegar su hechizo desde esa posición, sin atravesar los obstáculos.
\begin{itemize}
    \item El bosque se representa como una matriz booleana: 
    \begin{itemize}
        \item \texttt{false}: La celda está despejada (el hechizo puede pasar).  
        \item \texttt{true}: La celda está bloqueada (el hechizo no puede pasar).  
    \end{itemize}
    \item La entrada incluye la posición inicial del mago \((x, y)\).  
    \item El resultado debe ser una matriz que indique la distancia desde la posición inicial hasta cada celda accesible, respetando los obstáculos. Las celdas bloqueadas deben marcarse como \(-1\) en la salida.  
\end{itemize}

\subsection*{Ejemplo}
\text{Entrada:}
\begin{itemize}
        \item Posición inicial del mago: \((2, 2)\).  
        \item Bosque: 
        \[ 
        \begin{bmatrix}
            \texttt{false} & \texttt{false} & \texttt{false} & \texttt{false} & \texttt{false} \\
            \texttt{false} & \texttt{false} & \texttt{false} & \texttt{true}  & \texttt{false} \\
            \texttt{false} & \texttt{false} & \texttt{false} & \texttt{false} & \texttt{false} \\
            \texttt{false} & \texttt{false} & \texttt{false} & \texttt{true}  & \texttt{false} \\
            \texttt{false} & \texttt{false} & \texttt{false} & \texttt{false} & \texttt{false} \\
            \end{bmatrix}
        \]
\end{itemize}
\text{Salida:}
\[
\begin{bmatrix}
4 & 3 & 2 & 3 & 4  \\
3 & 2 & 1 & -1 & 3 \\
2 & 1 & 0 & 1  & 2 \\
3 & 2 & 1 & -1 & 3 \\
4 & 3 & 2 & 3  & 4 \\
\end{bmatrix}
\]

\section{Recorriendo islas}
El mapa de unas islas está modelado mediante un array bidimensional booleano, donde las casillas con valor true son secciones de tierra firme y las casillas con valor false son agua. Se desea encontrar cuál es la distancia mínima a recorrer entre el par de casillas \( (f_i, c_i) \) y \( (f_s, c_s) \) si solo se puede transitar por secciones de tierra firme.

\end{document}
